\documentclass[a4paper,11pt]{article}
\usepackage{exam-style-842-unit}
\examinehead{842 信息化概论 · 第09单元练习册}{练习二\quad 3S技术与农业专家系统}
\begin{document}

\unitcover{842 信息化概论 · 单元练习册}{第09单元\quad 信息化与农业应用}{练习二\quad 3S技术与农业专家系统（共 100 分）}{约 70 分钟}{遥感 RS 的原理与农业应用；地理信息系统 GIS 的组成、功能与农业应用；GPS／北斗的组成、定位原理与农业应用；3S 集成；农业专家系统的组成、工作原理与典型应用；农业数据库及管理系统；虚拟植物与虚拟农业。}

\begin{knowmap}
\textbf{3S 分工：}RS（面状·宏观·动态的信息源）· GIS（空间数据管理与分析平台）· GPS／北斗（点状·精确位置与导航）\\[2pt]
\textbf{3S 集成链：}RS 回答“哪里、发生了什么” $\rightarrow$ GIS 分析“为什么、怎么办” $\rightarrow$ GPS／北斗引导“精确到哪里作业”\\[2pt]
\textbf{遥感原理：}地物反射或辐射电磁波的光谱差异 $\rightarrow$ 传感器记录 $\rightarrow$ 图像处理与解译；分被动遥感与主动遥感\\[2pt]
\textbf{农业专家系统五部分：}知识库 · 推理机 · 综合数据库 · 人机接口 · 知识获取（解释机构可单列）\\[2pt]
\textbf{专家系统工作链：}事实输入 $\rightarrow$ 综合数据库 $\rightarrow$ 推理机匹配知识库规则 $\rightarrow$ 给出结论或建议 $\rightarrow$ 解释并更新知识\\[2pt]
\textbf{农业数据库与虚拟农业：}农业资源、品种、土壤、气象、病虫、市场等专题库由 DBMS 统一管理；虚拟植物以形态建模＋过程模拟＋可视化支撑虚拟农业。
\end{knowmap}

\begin{keyterms}
\textbf{遥感}：在不与目标接触的情况下，通过传感器接收地物反射或辐射的电磁波，经处理解译获取地物信息的技术。\\
\textbf{地理信息系统}：在计算机软硬件支持下，对具有空间位置的地理数据进行采集、存储、管理、分析、显示和应用的系统。\\
\textbf{全球定位系统}：由空间星座、地面控制和用户设备组成，为用户提供三维位置、速度和时间信息的卫星导航系统。\\
\textbf{北斗卫星导航系统}：我国自主研制的全球卫星导航系统，由空间段、地面段、用户段组成，兼具定位导航授时与短报文通信功能。\\
\textbf{3S 集成}：遥感、地理信息系统与全球定位系统在数据、平台与作业层面的综合集成，形成“监测—分析—作业”闭环。\\
\textbf{叠加分析}：把两个或多个空间图层按空间位置叠合，生成新属性与空间关系的 GIS 空间分析方法。\\
\textbf{归一化植被指数}：由近红外与红光波段反射率组合而成的植被指数（NDVI），用于反映植被覆盖与长势。\\
\textbf{农业专家系统}：把农业领域专家的知识与经验存入计算机，通过推理模拟专家思维、为农业生产提供决策建议的智能系统。\\
\textbf{知识库}：存放领域专家的知识、经验与规则，是专家系统推理的依据与核心。\\
\textbf{推理机}：依据综合数据库中的事实与知识库中的规则进行推理、得出结论的程序模块。\\
\textbf{综合数据库}：存放初始事实、用户输入数据和推理中间结果的动态数据库。\\
\textbf{知识获取}：把专家知识转化为计算机可表示的知识、并对知识库进行补充和维护的功能模块。\\
\textbf{虚拟植物}：以计算机模拟植物形态结构与生长发育过程、并可可视化表达的数字植物模型。\\
\textbf{虚拟农业}：利用虚拟现实、仿真与作物模型构建虚拟农场，对农业生产方案进行模拟、评价与培训的技术体系。
\end{keyterms}

\newpage
\begin{contrast}{遥感 / 地理信息系统 / 全球定位系统（每行依次为：获取的信息·核心技术·农业应用）}
\citem{遥感}{地物反射或辐射的电磁波信息，面状、宏观、动态}{传感器、遥感平台、图像处理与解译}{作物长势监测、面积调查、产量估测、灾害监测}
\citem{地理信息系统}{具有空间位置与拓扑关系的地理数据}{空间数据管理、空间查询与分析、可视化输出}{土壤与土地利用管理、适宜性评价、处方图制作}
\citem{全球定位系统}{精确的三维位置、速度与时间}{卫星测距定位、导航解算}{农机自动驾驶、精准播种、变量作业、农田测绘}
\end{contrast}

\begin{contrast}{知识库 / 综合数据库 / 推理机（每行依次为：存放内容·在系统中的作用·举例）}
\citem{知识库}{领域专家的知识、经验与规则}{推理的依据，是专家系统的核心}{“土壤 pH 低于 5.5 的酸性田块宜施石灰”}
\citem{综合数据库}{初始事实、用户输入与推理中间结果}{存放“当前问题”的动态数据}{某地块实测 pH 为 5.2、品种、前茬}
\citem{推理机}{推理策略与控制程序}{匹配事实与规则、推出结论}{正向推理得出“建议施石灰”}
\end{contrast}

\begin{contrast}{虚拟植物 / 虚拟农业（每行依次为：研究对象·关键技术·主要用途）}
\citem{虚拟植物}{植物个体或群体的形态结构与生长过程}{形态建模、生理生态过程模拟、可视化}{株型设计、栽培方案模拟、教学演示}
\citem{虚拟农业}{虚拟农场及其生产系统}{虚拟现实、仿真、作物模型与农业数据库集成}{生产方案优化、技术培训、园区规划}
\end{contrast}

\qtypename{一、单项选择题}{10 题，每题 3 分，共 30 分}
\q 遥感技术探测地物的基本原理是（\quad）
\choices{A. 利用地物反射或辐射电磁波的光谱差异}{B. 利用地物的重量差异}{C. 利用土壤酸碱度差异}{D. 利用农产品价格差异}

\q 在农业遥感中，常用来反映作物长势和植被覆盖状况的指标是（\quad）
\choices{A. 归一化植被指数（NDVI）}{B. 土壤容重}{C. 农药残留量}{D. 化肥价格指数}

\q 地理信息系统区别于一般数据库管理系统的最主要特点是（\quad）
\choices{A. 能管理具有空间位置和拓扑关系的地理数据}{B. 只能存储文字数据}{C. 只能绘制地图}{D. 不需要数据支撑}

\q 地理信息系统的核心功能是（\quad）
\choices{A. 空间数据的采集、存储、管理、分析与显示}{B. 文字排版与打印}{C. 视频剪辑与播放}{D. 语音识别与合成}

\q GPS 定位解算需要同时观测至少（\quad）颗卫星。
\choicesline{A. 2}{B. 3}{C. 4}{D. 6}

\q GPS 由三大部分组成，它们是（\quad）
\choices{A. 空间部分、地面控制部分、用户设备部分}{B. 感知层、网络层、应用层}{C. 硬件、软件、数据}{D. 输入、处理、输出}

\q 北斗卫星导航系统区别于 GPS 的特色功能是（\quad）
\choices{A. 短报文通信}{B. 只能定位不能导航}{C. 不能提供授时}{D. 只能用于军事领域}

\q 在 3S 集成中，承担空间数据管理与空间分析任务的是（\quad）
\choicesline{A. 遥感}{B. 地理信息系统}{C. 全球定位系统}{D. 虚拟现实}

\q 农业专家系统的核心是（\quad）
\choicesline{A. 知识库与推理机}{B. 显示器与键盘}{C. 网络设备}{D. 数据库软件}

\q 下列关于农业数据库及其管理系统的说法，正确的是（\quad）
\choices{A. 只能存储文本数据}{B. 用于农业数据的组织、存储、检索与维护}{C. 与农业生产实际无关}{D. 不需要数据库管理系统}

\qtypename{二、多项选择题}{4 题，每题 4 分，共 16 分；每题有两个或两个以上正确项}
\q 遥感技术在农业中的典型应用有（\quad）
\choices{A. 作物长势监测}{B. 种植面积调查}{C. 产量估测}{D. 直接完成施肥作业}

\q 地理信息系统的基本功能包括（\quad）
\choices{A. 数据采集与编辑}{B. 数据存储与管理}{C. 空间查询与分析}{D. 数据输出与可视化}

\q 农业专家系统一般包括（\quad）
\choicesline{A. 知识库}{B. 推理机}{C. 综合数据库}{D. 操作系统}

\q 北斗卫星导航系统在农业中的应用有（\quad）
\choices{A. 农机自动驾驶}{B. 精准播种}{C. 变量施肥}{D. 无人机植保导航}

\qtypename{三、判断题}{5 题，每题 2 分，共 10 分；正确的画“√”，错误的画“×”并写出改正后的表述}
\q 遥感只能被动接收地物反射的太阳光，因此夜间和阴天完全不能工作。\quad（\quad）\quad 改正：\fixblank

\q GIS 可以把土壤、地形、土地利用等图层叠加分析，用于土地适宜性评价。\quad（\quad）\quad 改正：\fixblank

\q GPS 定位只要接收到 3 颗卫星信号就能同时解算经度、纬度和高度。\quad（\quad）\quad 改正：\fixblank

\q 农业专家系统可以完全替代人类专家，不需要知识获取与知识更新。\quad（\quad）\quad 改正：\fixblank

\q 虚拟农业利用虚拟现实和仿真技术构建虚拟农场，可进行生产方案模拟与技术培训。\quad（\quad）\quad 改正：\fixblank

\qtypename{四、配对题}{6 对，每对 1 分，共 6 分；把 A 栏与 B 栏对应内容配对}
\q 将下列技术或部件与其主要功能配对，把 B 栏字母填入括号。
\par\vspace{0.2em}
\noindent\begin{tabular}{@{}p{0.42\textwidth}@{\hspace{1.0em}}p{0.46\textwidth}@{}}
（1）遥感 \hfill（\quad） & A. 空间数据管理与空间分析 \\
（2）地理信息系统 \hfill（\quad） & B. 地物光谱信息的获取 \\
（3）北斗卫星导航系统 \hfill（\quad） & C. 三维定位、导航与短报文通信 \\
（4）知识库 \hfill（\quad） & D. 存放领域专家的知识与规则 \\
（5）推理机 \hfill（\quad） & E. 依据事实与规则推出结论 \\
（6）虚拟植物 \hfill（\quad） & F. 植物形态与生长过程的可视化模拟 \\
\end{tabular}\par\vspace{0.3em}

\newpage
\qtypename{五、简答题}{4 题，每题 6 分，共 24 分}
\q 简述遥感技术的基本原理，并列举其在农业中的主要应用。\anslines{2.4cm}

\q 简述地理信息系统（GIS）的组成与基本功能，并举例说明其在农业中的应用。\anslines{2.4cm}

\q 简述 GPS／北斗卫星导航系统的组成与定位原理，并说明其在农业中的应用。\anslines{2.4cm}

\q 简述农业专家系统的组成与工作原理。\anslines{2.4cm}

\newpage
\qtypename{六、综合题}{1 题，共 14 分}
\q 某市建设“天空地”一体化农情监测与精准作业服务体系：卫星遥感定期获取全市作物分布与长势，无人机对重点地块开展病虫害巡查，地面物联网采集土壤墒情，GIS 平台叠加土壤、气象、产量图层划分管理单元并生成处方图，农业专家系统给出施肥与防治建议，北斗终端引导农机实施变量作业。请回答：（1）遥感、GIS、GPS／北斗在该体系中各自承担什么任务？三者如何集成为“监测—分析—作业”链条？（2）农业专家系统嵌入决策环节后，其工作过程是怎样的？（3）说明该体系对农业生产的作用与价值。\anslines{6cm}

\end{document}
